\documentclass[11pt]{article}

\usepackage[margin=1in]{geometry}
\usepackage{amsmath, amssymb}
\usepackage{listings}
\usepackage{xcolor}
\usepackage{hyperref}
\usepackage{enumitem}
\usepackage{titlesec}
\usepackage{parskip}
\usepackage{booktabs}

% Code listing style
\lstset{
    basicstyle=\ttfamily\small,
    keywordstyle=\color{blue}\bfseries,
    commentstyle=\color{gray}\itshape,
    stringstyle=\color{red},
    frame=single,
    breaklines=true,
    showstringspaces=false,
    tabsize=4
}

\title{\textbf{CS3210: Operating Systems} \\ Lecture 10 --- Interrupts and Concurrency (Part 1)}
\author{John Kim \\ Georgia Institute of Technology}
\date{February 17, 2026}

\begin{document}

\maketitle
\tableofcontents
\newpage

% ============================================================
\section{Swapping Wrap-Up: Page Eviction Policies}
% ============================================================

\subsection{The Problem of Choosing Pages to Evict}

\textbf{Swapping} is the OS mechanism for saving memory pages to disk when physical RAM is exhausted, then reloading them on demand. The central challenge is selecting \emph{which} pages to evict. The \textbf{ideal} candidate is the page that will be accessed \emph{latest} in the future---this is the \textbf{Belady's optimal} policy. However, this is impossible to compute in practice because future access patterns are unknown.

The standard approximation is \textbf{Least Recently Used (LRU)}: evict the page that has not been accessed for the longest time. The intuition is that recently-used pages are more likely to be needed again soon (temporal locality). The difficulty is that precise LRU requires tracking a timestamp on every memory access---including every load and store---which imposes prohibitive overhead even when memory is plentiful.

\subsection{Tracking LRU: Complexity Analysis}

To maintain exact LRU, the OS would need to update a data structure on each memory access. The cost of maintaining this structure ranges from $O(n)$ down to $O(1)$ with the right design, but the fundamental problem is that \emph{every} memory access---including reads---would require trapping into the kernel to update the timestamp. This means a page fault on every load and store, which is completely impractical.

\subsection{The Clock Algorithm}

The \textbf{clock algorithm} is a hardware-assisted LRU approximation that avoids per-access traps. Every page table entry (PTE) contains an \textbf{access bit} $a$. The hardware sets $a = 1$ whenever a page is read or written. The OS maintains pages in a \emph{circular list} and advances a ``hand'' pointer when eviction is needed.

The eviction procedure works as follows:
\begin{enumerate}
    \item Examine the page at the current hand position.
    \item If $a = 1$, clear the bit ($a \leftarrow 0$) and advance the hand. This gives the page a second chance.
    \item If $a = 0$, evict this page. It has not been accessed since the last time the hand passed.
\end{enumerate}

The key insight connecting hardware and software: when $a = 0$, the OS also clears \texttt{PTE\_P} (the present bit). This causes the next access to the page to generate a \textbf{page fault}, at which point the OS sets $a = 1$ and \texttt{PTE\_P} again. This way, trapping happens only when the access bit needs to be updated from $0$ to $1$---not on every access.

\subsubsection*{Worst-Case Analysis}

If \emph{all} pages have $a = 1$ when eviction begins, the clock hand must sweep the entire circular list, clearing bits one by one, before finding any page to evict. In the worst case this is $O(n)$ where $n$ is the number of pages---but it degenerates only when every page is actively used, exactly the situation where no good eviction candidate exists regardless of algorithm.

\begin{center}
\begin{tabular}{ll}
\toprule
\textbf{Policy} & \textbf{Characteristic} \\
\midrule
Optimal (Belady) & Evict page accessed latest in the future; impractical \\
LRU (exact) & Evict least-recently used; requires per-access trap \\
Clock (LRU approx.) & Circular scan with access bit; hardware-assisted, practical \\
Random & Simple, unpredictable; baseline comparison \\
\bottomrule
\end{tabular}
\end{center}

% ============================================================
\section{Interrupts, Exceptions, and System Calls}
% ============================================================

\subsection{Why Interrupts?}

Without interrupts, the only alternative is \textbf{polling}: the CPU repeatedly checks each I/O device to see if it needs attention. For \emph{infrequent} events (e.g., a keystroke) this wastes enormous CPU cycles. For \emph{frequent} events polling may not even be fast enough to keep up. Interrupts solve both problems: the device signals the CPU only when action is required, and the CPU handles it immediately.

Events requiring interrupt handling fall into several categories:
\begin{itemize}
    \item \textbf{Exceptions} (synchronous): page faults (vector 14 in x86), arithmetic errors, protection faults.
    \item \textbf{Hardware interrupts} (asynchronous): timer ticks (vector 32 $\to$ scheduler), disk completion, network packet arrival.
    \item \textbf{Software interrupts}: the \texttt{INT} instruction, used for system calls in xv6 (\texttt{INT 0x40}).
    \item \textbf{Non-Maskable Interrupts (NMI)}: power failure, memory corruption; cannot be suppressed.
    \item \textbf{Debugging}: software breakpoints.
    \item \textbf{Communication}: Inter-processor interrupts (IPI) on multiprocessor systems.
\end{itemize}

Every interrupt and exception type is assigned a fixed \textbf{vector number} (0--255). When an interrupt fires, the hardware uses the vector to index into the \textbf{Interrupt Descriptor Table (IDT)} and jump to the registered handler.

\subsection{Types of Interrupts in x86}

\subsubsection*{Non-Maskable Interrupts (NMI)}

\textbf{NMI} (vector 2) is \emph{never} ignored by the CPU, regardless of the interrupt-enable flag. It signals catastrophic hardware events such as memory ECC errors or power failure. In xv6, receiving an NMI causes the kernel to halt---there is no safe recovery path.

\subsubsection*{Maskable Interrupts (INTR)}

\textbf{Maskable interrupts} are suppressed when the \textbf{IF} (interrupt-enable) flag in \texttt{EFLAGS} is cleared. The x86 instructions \texttt{sti} (set interrupt flag) and \texttt{cli} (clear interrupt flag) enable and disable maskable interrupts, respectively. The CPU sends an \textbf{INTA} (interrupt acknowledge) signal back to the hardware controller after handling each maskable interrupt.

The \textbf{Programmable Interrupt Controller (PIC 8259A)} aggregates multiple hardware interrupt lines into the single \texttt{INTR} pin of the CPU. It is responsible for encoding which device triggered the interrupt and presenting the correct vector number.

\subsubsection*{Software Interrupts (INT)}

The \texttt{INT n} instruction deliberately triggers the interrupt handler for vector $n$. It is the x86 mechanism for \textbf{system calls}: xv6 uses \texttt{INT 0x40} to enter the kernel. The complementary instruction \texttt{iret} (interrupt return) restores the saved processor state and returns to user space.

\subsection{The INT Instruction: Step-by-Step}

When \texttt{INT 0x40} executes in user space, the CPU performs the following sequence automatically:
\begin{enumerate}
    \item Determine the vector number (0x40).
    \item Fetch the 8-byte \textbf{IDT descriptor} for that vector from the table whose base address is in the \texttt{IDTR} register.
    \item Check that the current privilege level (CPL) $\leq$ descriptor privilege level (DPL)---this prevents user code from directly invoking kernel-only handlers.
    \item If a privilege-level change is required, save \texttt{SS} and \texttt{ESP} internally, then load the kernel \texttt{SS} and \texttt{ESP} from the \textbf{Task State Segment (TSS)}. This switches to the kernel stack.
    \item Push (in order): user \texttt{SS}, user \texttt{ESP}, user \texttt{EFLAGS}, user \texttt{CS}, user \texttt{EIP}. These form the bottom of the \emph{trap frame}.
    \item Clear certain \texttt{EFLAGS} bits (e.g., the interrupt-enable flag is cleared to prevent re-entrant interrupts during handler setup).
    \item Set \texttt{CS} and \texttt{EIP} from the IDT descriptor, transferring control to the handler.
\end{enumerate}

The net effect is that the hardware automatically saves the minimal state needed to return to user space, and switches to the kernel stack before invoking the handler.

% ============================================================
\section{Interrupt Handling in xv6}
% ============================================================

\subsection{High-Level Flow}

Handling any interrupt in xv6 follows this general pattern:
\begin{enumerate}
    \item Identify the interrupt vector.
    \item Verify the interrupt is permitted at the current privilege level.
    \item Switch from the user stack to the kernel (handler) stack via the TSS.
    \item Save all user-space registers (building the full \emph{trap frame}).
    \item Jump to the C trap handler.
    \item On return, restore the saved registers and execute \texttt{iret} to resume user space.
\end{enumerate}

\subsection{The Trap Frame (\texttt{x86.h})}

The \textbf{trap frame} is the data structure that captures the complete CPU state at the moment of an interrupt. In xv6, it is defined in \texttt{x86.h}. The lower portion (pushed first) is filled by the hardware during the \texttt{INT} instruction; the upper portion is filled by software in \texttt{trapasm.S}. Together they hold all general-purpose registers, segment registers, and the hardware-pushed \texttt{EIP}/\texttt{CS}/\texttt{EFLAGS}/\texttt{ESP}/\texttt{SS}, providing everything needed to restore user state on \texttt{iret}.

\subsection{xv6 Interrupt Source Files}

The interrupt infrastructure in xv6 spans several files, each with a distinct responsibility:

\begin{itemize}
    \item \textbf{\texttt{usys.S}}: Defines the user-space system call stubs. Each stub executes \texttt{INT 0x40} to enter the kernel.
    \item \textbf{\texttt{trapasm.S}}: Contains \texttt{alltraps}, the unified entry point for all interrupts. It pushes the remaining registers to complete the trap frame, then calls the C function \texttt{trap()}. It also contains \texttt{trapret}, which restores state and executes \texttt{iret} to return to user space.
    \item \textbf{\texttt{vector.S}}: Auto-generated file defining 256 short stubs---one per interrupt vector. Each stub pushes the vector number (and a dummy error code for vectors that do not push one automatically), then jumps to \texttt{alltraps}.
    \item \textbf{\texttt{trap.c}}: The main C-level trap handler. This file contains the \texttt{trap()} function that dispatches on the vector number, handling system calls, hardware interrupts, and exceptions. It is the file most relevant to Lab 2 and Lab 3.
\end{itemize}

\subsection{The Interrupt Descriptor Table (IDT)}

The \textbf{IDT} is an array of 256 eight-byte entries analogous to the GDT. Each entry is called a \emph{gate descriptor} and specifies:
\begin{itemize}
    \item The handler's segment selector and offset (the address of the handler code).
    \item The privilege level required to invoke this gate via a software \texttt{INT} instruction.
    \item The gate type (interrupt gate vs. trap gate; interrupt gates clear \texttt{IF} on entry).
\end{itemize}

The CPU locates the IDT via the \textbf{IDTR} register, which holds the base address and size of the table. The \texttt{lidt} instruction loads \texttt{IDTR}. In xv6, the IDT is initialized at boot in \texttt{trap.c} using a helper \texttt{SETGATE} macro that populates each descriptor.

% ============================================================
\section{Introduction to Concurrency}
% ============================================================

\subsection{Interrupts as Pseudo-Concurrency}

Even on a single CPU, interrupts create a form of \textbf{pseudo-concurrency}: an interrupt handler can run at any point between any two instructions of the main kernel code. If both the interrupted kernel code and the interrupt handler access shared data, a \textbf{race condition} arises. This is logically equivalent to the concurrency problems that occur with multiple threads or processors.

\subsection{The Deadlock Problem with Locks and Interrupts}

Suppose the kernel acquires a spinlock protecting shared data, and then an interrupt fires before the lock is released:

\begin{enumerate}
    \item Kernel code calls \texttt{Lock()} and acquires the lock.
    \item Interrupt arrives; the CPU jumps to the interrupt handler.
    \item The handler also calls \texttt{Lock()} on the \emph{same} lock.
    \item Since the lock is already held, the handler spins indefinitely---\textbf{deadlock}.
\end{enumerate}

The handler cannot be made to ignore locks (that would break mutual exclusion). The solution is to \textbf{delay the interrupt}: whenever a spinlock is held, interrupts must be disabled. This is the approach xv6 takes---\texttt{acquire()} calls \texttt{cli} to disable interrupts, and \texttt{release()} calls \texttt{sti} to re-enable them only after the lock is released.

In xv6, the lock primitives therefore have the following structure conceptually:

\begin{lstlisting}[language=C, caption={Interrupt-safe acquire/release in xv6 (conceptual)}]
acquire(lock):
    interrupt_disable()   // pushcli in xv6
    spinlock_acquire(lock)

release(lock):
    spinlock_release(lock)
    interrupt_enable()    // popcli in xv6
\end{lstlisting}

\textbf{Mutual exclusion with an interrupt handler is achieved through interrupt disable.} This is a key design principle in xv6: the combination of spinlocks and \texttt{cli}/\texttt{sti} prevents the interrupt-deadlock scenario entirely.

% ============================================================
\section{Summary}
% ============================================================

\begin{itemize}
    \item \textbf{Clock algorithm}: An LRU approximation using the hardware access bit ($a$) and a circular page list; avoids per-access traps by piggybacking on page faults.
    \item \textbf{Interrupt vectors}: Every exception and interrupt maps to a vector number (0--255); the IDT maps vectors to kernel handler addresses.
    \item \textbf{NMI vs.\ INTR vs.\ INT}: Non-maskable interrupts cannot be suppressed; maskable interrupts are controlled by the \texttt{IF} flag (\texttt{sti}/\texttt{cli}); software interrupts (\texttt{INT n}) are used for system calls.
    \item \textbf{INT instruction mechanics}: Automatically switches to the kernel stack (via TSS), pushes user state onto it, and transfers control to the IDT handler.
    \item \textbf{xv6 interrupt path}: \texttt{vector.S} $\to$ \texttt{alltraps} (trapasm.S) $\to$ \texttt{trap()} (trap.c) $\to$ \texttt{trapret} $\to$ \texttt{iret}.
    \item \textbf{IDT}: 256-entry table of gate descriptors; base address loaded into \texttt{IDTR} via \texttt{lidt}; initialized at boot in \texttt{trap.c}.
    \item \textbf{Pseudo-concurrency}: Interrupts create race conditions on a single CPU; xv6 prevents interrupt-induced deadlock by disabling interrupts while holding spinlocks.
\end{itemize}

% ============================================================
\section{References}
% ============================================================

\begin{itemize}
    \item Cox, R.; Kaashoek, F.; Morris, R. ``xv6: a simple, Unix-like teaching operating system.'' \textit{MIT CSAIL}. \url{https://pdos.csail.mit.edu/6.828/2018/xv6/book-rev11.pdf}
    \item Intel. ``Intel 64 and IA-32 Architectures Software Developer's Manual, Volume 3A: System Programming Guide.'' Intel Corporation.
    \item Arpaci-Dusseau, R. H.; Arpaci-Dusseau, A. C. ``Beyond Physical Memory: Mechanisms.'' \textit{Operating Systems: Three Easy Pieces}. Arpaci-Dusseau Books. \url{https://pages.cs.wisc.edu/~remzi/OSTEP/}
    \item TutorialsPoint. ``Microprocessor 8086 Interrupts.'' \url{https://www.tutorialspoint.com/microprocessor/microprocessor_8086_interrupts.htm}
\end{itemize}

% ============================================================
\section*{Personal Notes}
% ============================================================
% Add your own notes below this line

\end{document}
